\documentclass[12pt,a4paper]{article}
\usepackage[utf8]{inputenc}
\usepackage{amsmath}
\usepackage{amsfonts}
\usepackage{amssymb}
\usepackage{tikz}
\usepackage{geometry}
\geometry{margin=1in}

\title{Lattice Path from (0,0) to ({{ m }}, {{ n }})}
\author{MCP Adapter LaTeX Server}
\date{\today}

\begin{document}

\maketitle

\section{Lattice Path Visualization}

This document shows a lattice path from the origin $(0,0)$ to the point $({{ m }}, {{ n }})$ on a rectangular grid.

\begin{center}
\begin{tikzpicture}[scale=0.8]
    % Define the grid size with some padding
    \def\gridwidth{{{ m + 2 }}}
    \def\gridheight{{{ n + 2 }}}
    
    % Draw the grid
    \foreach \x in {0,...,\gridwidth}
        \draw[very thin,gray] (\x,0) -- (\x,\gridheight);
    \foreach \y in {0,...,\gridheight}
        \draw[very thin,gray] (0,\y) -- (\gridwidth,\y);
    
    % Draw axes with labels
    \draw[thick] (0,0) -- (\gridwidth,0);
    \draw[thick] (0,0) -- (0,\gridheight);
    
    % Label axes
    \foreach \x in {0,1,...,{{ m }}}
        \node[below] at (\x,0) {\x};
    \foreach \y in {0,1,...,{{ n }}}
        \node[left] at (0,\y) {\y};
    
    % Highlight the target rectangle
    \draw[thick,red] (0,0) rectangle ({{ m }},{{ n }});
    
    % Draw a sample lattice path (monotonic path)
    \draw[very thick,blue,->] (0,0) 
    {% for i in range(m) %}
        -- ({{ i + 1 }},0)
    {% endfor %}
    {% for j in range(n) %}
        -- ({{ m }},{{ j + 1 }})
    {% endfor %};
    
    % Mark start and end points
    \fill[green] (0,0) circle (0.1) node[below left] {Start (0,0)};
    \fill[red] ({{ m }},{{ n }}) circle (0.1) node[above right] {End ({{ m }},{{ n }})};
    
    % Add grid points
    \foreach \x in {0,...,{{ m }}}
        \foreach \y in {0,...,{{ n }}}
            \fill[black] (\x,\y) circle (0.05);
\end{tikzpicture}
\end{center}

\section{Mathematical Properties}

\subsection{Path Description}
\begin{itemize}
    \item \textbf{Start Point:} $(0, 0)$
    \item \textbf{End Point:} $({{ m }}, {{ n }})$
    \item \textbf{Horizontal Steps:} {{ m }}
    \item \textbf{Vertical Steps:} {{ n }}
    \item \textbf{Total Steps:} {{ m + n }}
\end{itemize}

\subsection{Combinatorial Analysis}
The number of monotonic lattice paths from $(0,0)$ to $({{ m }}, {{ n }})$ is given by the binomial coefficient:

\[
\binom{{{ m + n }}}{{{ m }}} = \binom{{{ m + n }}}{{{ n }}} = \frac{({{ m + n }})!}{{{ m }}! \cdot {{ n }}!}
\]

This represents the number of ways to choose {{ m }} horizontal steps out of {{ m + n }} total steps.

\subsection{Path Enumeration}
Each path can be encoded as a sequence of {{ m + n }} moves, where:
\begin{itemize}
    \item R represents a right move (horizontal step)
    \item U represents an up move (vertical step)
\end{itemize}

The path shown above corresponds to the sequence: $\underbrace{\text{R R} \ldots \text{R}}_{{{ m }} \text{ times}} \underbrace{\text{U U} \ldots \text{U}}_{{{ n }} \text{ times}}$

\subsection{Applications}
Lattice paths have applications in:
\begin{itemize}
    \item Combinatorics and counting problems
    \item Probability theory (random walks)
    \item Computer science (dynamic programming)
    \item Physics (statistical mechanics)
\end{itemize}

\end{document}